% ch1_answers.tex -- answers to the Quick check questions of Module 8, Chapter 1.
\begin{chapteranswers}{1}{Origins of the Elements: the Big Bang and the Stars}
  \item Quarks $\to$ protons $\to$ helium nuclei $\to$ atoms $\to$ first stars.
        (Quarks join into protons at about \qty{e-6}{s}; helium nuclei form at 3 to 20
        minutes; atoms at \num{380000} years; stars after about 200 million years.)
  \item Two particles of mass $m_\mathrm{e}$ are made, so
        $E = 2m_\mathrm{e}c^2 = 2 \times \qty{9.109e-31}{kg} \times (\qty{3.00e8}{m.s^{-1}})^2
        = \qty{1.6396e-13}{J}$.
        In eV: $\dfrac{\qty{1.6396e-13}{J}}{\qty{1.602e-19}{J.eV^{-1}}} = \qty{1.0235e6}{eV}$.
        Answer: \final{$E = \qty{1.02}{MeV}$}.
  \item A proton pair needs at least $2m_\mathrm{p}c^2 \approx \qty{1880}{MeV}$, but an
        electron pair needs only about \qty{1.02}{MeV}. As the Universe expanded and
        cooled, the photons' energies fell. They fell below \qty{1880}{MeV} first, while
        the Universe was still extremely hot, so proton pairs stopped first. Photons
        could still make electron pairs until their energies fell below about
        \qty{1}{MeV}, a few seconds after the Big Bang.
  \item $\Delta\lambda = \qty{397.3}{nm} - \qty{393.4}{nm} = \qty{3.9}{nm}$. Then
        $v = \dfrac{c\,\Delta\lambda}{\lambda_0}
        = \dfrac{\qty{3.00e8}{m.s^{-1}} \times \qty{3.9}{nm}}{\qty{393.4}{nm}}
        = \qty{2.974e6}{m.s^{-1}} = \qty{2974}{km.s^{-1}}$ (the nm cancel).
        $d = \dfrac{v}{H_0} = \dfrac{\qty{2974}{km.s^{-1}}}{\qty{70}{km.s^{-1}.Mpc^{-1}}}
        = \qty{42.49}{Mpc}$.
        Answers: \final{$v \approx \qty{3.0e3}{km.s^{-1}}$} and \final{$d \approx \qty{42}{Mpc}$}
        (2 significant figures, because $\Delta\lambda$ has only 2).
  \item She also sees almost every distant galaxy redshifted, moving away from
        \emph{her}, with speed proportional to distance. The Milky Way moves away from
        her at $v = H_0 d = \qty{70}{km.s^{-1}.Mpc^{-1}} \times \qty{300}{Mpc}
        = \qty{2.1e4}{km.s^{-1}}$. Her view is the same as ours, as the cosmological
        principle says. This happens because space itself expands everywhere, like
        raisins in rising bread. So no galaxy is at the centre: the Universe has no
        centre.
  \item The CMB is light from the hot early Universe. It was set free about
        \num{380000} years after the Big Bang, when atoms formed and the Universe was
        about \qty{3000}{K}. Then it was mostly infrared and visible light (peak about
        \qty{966}{nm}). As the light travelled, the expansion of space stretched its
        wavelengths about 1100 times, to about \qty{1}{mm}: microwaves. Dust sits in
        particular places, but the CMB is the same in every direction, so it fills all
        of space. Its microwave glow is what an expanding, cooling Universe predicts.
  \item Both lines come from the same gap between two energy levels of sodium atoms.
        For one photon of \qty{589}{nm} (as in \cref{ch1:we:lines}):
        $f = \dfrac{c}{\lambda} = \qty{5.09e14}{Hz}$, so
        $E = hf = \qty{3.37e-19}{J} = \qty{2.11}{eV}$. This equals the gap.\newline
        \textbf{Lamp (emission):} an electric current passes through hot, low-pressure
        sodium gas. Collisions lift electrons to the higher level. Each electron then
        drops back and gives out one photon of exactly \qty{2.11}{eV}. So we see a
        bright line on a dark background.\newline
        \textbf{Sunlight (absorption):} the Sun's hot, dense inside gives a continuous
        spectrum. Cooler sodium gas in its outer layers absorbs photons of exactly
        \qty{2.11}{eV}, which lift electrons up. The energy is given out again in all
        directions, so less of this wavelength reaches us: a dark line.\newline
        The gap is the same in both, so both lines are at \qty{589}{nm}.
  \item $\lambda_\mathrm{max} = \dfrac{b}{T} = \dfrac{\qty{2.898e-3}{m.K}}{\qty{25000}{K}}
        = \qty{1.16e-7}{m} = \qty{116}{nm}$.\\[3pt] Answer: \final{\qty{116}{nm}}, in the
        ultraviolet. A star at \qty{25000}{K} is class B (it looks blue-white).
  \item Visible hydrogen lines are absorbed only by electrons that start in level
        $n = 2$. Star P (class O9) is very hot, so most of its hydrogen is ionised: the
        electron has been removed, so it cannot absorb. Star Q (class M3) is cool, so
        almost all its hydrogen electrons are in $n = 1$, and few can absorb these
        wavelengths. Only in between (class A, about \qty{10000}{K}) is the share of
        electrons in $n = 2$ largest.
  \item Lines at the same wavelengths mean the same elements (and similar
        temperatures). Broad lines come from a dense, high-pressure atmosphere, as in a
        small main sequence star. Narrow lines come from a thin, low-density atmosphere.
        So \final{star Y}, with the narrow lines, is more likely to be a giant.
  \item At the top right. These stars are cool, so each square metre gives out
        little light. To be very luminous they must have a huge surface: they are
        giants or supergiants.
  \item Star A: a smaller magnitude means more luminous. The difference is
        $10 - (-5) = 15$ magnitudes, which is three steps of 5 magnitudes. Each step
        is a factor of 100, so A is about $100^3 = 10^6$ times more luminous.
        Answer: \final{star A, about a million times}.
  \item A $10\,M_\odot$ star has stronger gravity, so its core is much hotter.
        Fusion runs very much faster. It has 10 times more fuel, but it uses it
        thousands of times faster. So it runs out of core hydrogen after tens of
        millions of years, not 10 billion years.
  \item $\Delta m = 3 \times \qty{4.002603}{u} - \qty{12.000000}{u} = \qty{0.007809}{u}$.
        $E = \num{0.007809} \times \qty{931.5}{MeV} = \qty{7.274}{MeV}$.
        Answer: \final{$E = \qty{7.27}{MeV}$}.
  \item $A$: $13 + 1 = 14$. $Z$: $6 + 1 = 7$, which is nitrogen. Answer:
        \final{$\chOneNuc{14}{7}{N}$} (step 3 of the CNO cycle).
  \item Any two of these:\newline
        $\bullet$ The p--p chain is the main route in the Sun and smaller stars. The CNO
        cycle is the main route in stars above about $1.3\,M_\odot$.\newline
        $\bullet$ The CNO cycle needs a hotter core (above about 18 million K).\newline
        $\bullet$ The CNO cycle needs carbon, nitrogen or oxygen already in the star, as
        catalysts. The p--p chain needs only hydrogen.
  \item Fusion releases energy only when the new nucleus is more tightly bound than the
        old ones. This means its protons and neutrons hold together more strongly, so
        some mass is lost as energy. This is true only up to iron. Fusing iron into heavier nuclei takes
        in energy instead of releasing it, so it cannot keep the core hot. Heavier
        elements are made instead by neutron capture, in supernovae, neutron-star
        collisions and old giant stars.
\end{chapteranswers}
